\documentclass[10pt,a4paper]{amsart}
\usepackage[margin=19mm]{geometry}
\usepackage{amsmath,amssymb,mathtools}
\usepackage[scr=rsfs,cal=euler]{mathalfa}
\usepackage{xfrac}
\usepackage[unicode,hidelinks,bookmarks=false]{hyperref}
\newcommand{\ie}{i.e.\ }
\newcommand{\cf}{cf.\ }
\newcommand{\resp}{respectively\ }
\newcommand{\Subsection}{\subsection}
\providecommand{\nobreakdash}{\nobreak\hbox{-}}
\setlength{\emergencystretch}{2em}
\multlinegap0pt
\usepackage{bm}
\usepackage{bbm}
\usepackage{dsfont}
\usepackage{xspace}
\usepackage{cancel}
\usepackage{mathtools}
\usepackage{amsthm}
\makeatletter
\@ifundefined{lemma}{\newtheorem{lemma}{Lemma}}{}
\@ifundefined{proposition}{\newtheorem{proposition}[lemma]{Proposition}}{}
\makeatother
\newcommand{\bb}[1]{{\mathbb #1}} 
\newcommand{\f}{{\mf f}}
\newcommand{\mf}[1]{{\mathfrak #1}} 
\title[A Central Limit Theorem for intransitive dice]{A Central Limit Theorem for intransitive dice}
\author{Luis G. Coelho, Tertuliano Franco, Lael V. Lima, Jo\~ao P. C. de Paula, Jo\~ao V. A. Pimenta, Guilherme L. F. Silva, Daniel Ungaretti}
\date{Source: arXiv:2310.17083v2 (CC BY 4.0)}
\begin{document}
\maketitle

\noindent\emph{Source and licence.} A Central Limit Theorem for intransitive dice. Version arXiv:2310.17083v2 (CC BY 4.0), distributed under the Creative Commons Attribution 4.0 International licence, as recorded in the arXiv licence metadata for this exact version and shown on the abstract page https://arxiv.org/abs/2310.17083v2. Full rights notice: licenses/arxiv-2310.17083-CC-BY-4.0.txt.\par\smallskip

\noindent\textit{Editorial scope.} Counted unit: the Proposition whose source label is \texttt{prop:h2\_eigenspace\_zero} in the article (source0.tex lines 2002--2005, source label \texttt{prop:h2\_eigenspace\_zero}), delivered together with the article's own complete original proof of it (source0.tex lines 2006--2130, one \texttt{proof} environment of 125 lines). No mathematics is written, repaired or re-derived: statement and proof are the article's own text line for line. Required context retained uncounted: eq:covariancematrixCLT (lines 583--593); eq:gammakfrakf (lines 1840--1842); eq:def\_event\_Ak (lines 1958--1961); eq:prob\_cup\_Bk\_2 (lines 1989--1995). Every definition, display and label of those lines is retained unchanged. Omitted from this package and not redistributed: 1--582, 594--1839, 1843--1957, 1962--1988, 1996--2001, 2131--3085. Nothing inside a retained excerpt is omitted or altered: every retained body is delivered exactly as the article's lines are written, so no omission receipt of any kind accompanies this package. Citations: the retained excerpts carry no citation command at all, so no bibliography entry of the article is transcribed and no reference is redistributed. Reference adaptation: none was needed and none was made; every reference inside the delivered unit resolves inside it, so no unresolved reference or citation is produced. Printed slips kept literal: no printed word, symbol, hypothesis or number of the counted statement or of its proof was altered. Layout and class: the article's own class is \texttt{amsart} with packages including inputenc, babel, helvet, courier, upgreek, amsthm, bm, bbm, float, amsmath, amssymb, tikz, enumerate, dsfont; the wrapper is the lane's pinned \texttt{amsart} editorial wrapper at 10pt on A4, which restates the theorem-like environments the retained text needs and adds the article's own macro definitions that the retained text needs (\texttt{\textbackslash bb}, \texttt{\textbackslash f}, \texttt{\textbackslash mf}), with extra wrapper packages (bm, bbm, dsfont, xspace, cancel, mathtools). The delivered numbering is the unit's own. Assets: no retained span contains a figure environment, a graphics-inclusion command, a PDF-page inclusion, a table or a TikZ picture; no image file is copied into this package, the package declares no asset dependency and no image file is redistributed with it. Licence: version arXiv:2310.17083v2 of this article is distributed under the Creative Commons Attribution 4.0 International licence, as recorded in the arXiv licence metadata for this exact version and shown on the abstract page https://arxiv.org/abs/2310.17083v2. A full rights notice is delivered at licenses/arxiv-2310.17083-CC-BY-4.0.txt. Characters: every retained excerpt is pure ASCII, so the delivered engine, XeLaTeX with its default Latin Modern text font, reports no missing character.
\begin{equation}\label{eq:covariancematrixCLT}
\Sigma \;=\;\left(
\begin{array}{cccccccc}
1 & \gamma_2(\infty) & 0 & \cdots  & 0 &\gamma_1(\infty)\\
\gamma_2(\infty) & 1 & \gamma_3(\infty) & \cdots & 0 & 0\\
0 & \gamma_3(\infty) & 1 & \cdots& 0 & 0\\
\vdots & \vdots & \vdots & \ddots & \vdots & \vdots\\
0 & 0 & 0 & \cdots & 1 & \gamma_{\ell}(\infty)\\
\gamma_1(\infty) & 0 & 0 & \cdots & \gamma_{\ell}(\infty) & 1
\end{array}\right),
\end{equation}

\begin{equation}\label{eq:gammakfrakf}
\gamma_k(\infty)\;=\;-\frac{\f_{k-1}\f_k\mf f_{k+1}}{\sqrt{\f_{k-1}\mf f_k(\f_{k-1}+\f_{k})}\sqrt{\f_{k}\f_{k+1}(\f_{k}+\f_{k+1})}}\,,
\end{equation}

\begin{equation}
\label{eq:def_event_Ak}
    A_k \;\coloneqq\; \{U_k \le a_k\}\,.
\end{equation}

\begin{align}
\bb P(\cup B_k)
    &\;=\; 1 - \bb P(\cap B^c_k)
    \;=\; 1 - \bb P(\cap A^c_k) -  \bb P(\cap A_k)\nonumber\\
    &\;=\; 1 - a_1 \cdots a_\ell - (1 - a_1) \cdots (1 - a_\ell) \nonumber\\
    &\;=\; 1 - 2 \prod \frac{\f_k}{\f_k + \f_{k+1}}     = 1 - 2 (-1)^\ell \prod \gamma_k\,.\label{eq:prob_cup_Bk_2}
\end{align}

\begin{proposition}
\label{prop:h2_eigenspace_zero}
Let $\Sigma$ be as in \eqref{eq:covariancematrixCLT} with coefficients $\gamma_k=\gamma_k(\infty)$ as in \eqref{eq:gammakfrakf}. Then $0$ is an eigenvalue of $\Sigma$, its eigenspace has dimension 1 and is generated by a vector $x \in (0,\infty)^{\ell}$.
\end{proposition}

\begin{proof}
Let $x=(x_1, \dots, x_\ell)$ be a non-zero vector satisfying $\Sigma x = 0$.
Then, for every $k \in [\ell]$ we have
\begin{equation}
\label{eq:eigen_0_Lk} \tag{$L_k$}
\gamma_{k} x_{k-1} + x_{k}+\gamma_{k+1} x_{k+1} \;=\;0\,.
\end{equation}

It is possible to solve the system of equations above explicitly, but the formulas obtained this way are cumbersome. Instead, we show that if some coordinate $x_j$ is positive then $x_{j+1}$ is positive as well. Since we can always choose the sign of one entry of $x$, by the cyclic symmetry of the problem we then conclude that there is $x \in (0,\infty)^{\ell}$ with $\Sigma x = 0$, as wanted.
Therefore, assume without loss of generality that $x_{\ell-1}\geq 0$. From
$(L_1)-\gamma_1(L_\ell)$, we obtain
\begin{equation*}
 -\gamma_1\gamma_\ell x_{\ell-1} +(1-\gamma_1^2)x_1+\gamma_2 x_2\;=\;0\,.
\end{equation*}
Defining $P_0 := 1$ and $P_1 := 1 - \gamma_1^2$, the equation above becomes
\begin{equation}
 -\gamma_1\gamma_\ell x_{\ell-1} +P_1 x_1+\gamma_2 P_0 x_2\;=\;0 \tag{$L'_1$}\,.
\end{equation}
Equation $(L'_1)$ relates $x_{\ell-1}$ to $x_1$ and $x_2$. By successive applications
of the same reasoning we can relate $x_{\ell-1}$ to $x_k$ and $x_{k+1}$ for any
$k$. Indeed, suppose that it holds
\begin{equation}
\tag{$L'_k$}
(-1)^k\gamma_1\dots\gamma_k\gamma_\ell x_{\ell-1}+P_k x_k+\gamma_{k+1}
    P_{k-1}x_{k+1}=0,
\end{equation}
for some already defined $P_{k-1}$ and $P_k$.
Then, from $P_k(L_{k+1})-\gamma_{k+1}(L'_k)$, we obtain
\begin{align*}
    &\bigl(P_k \gamma_{k+1} x_{k}
    + P_k x_{k+1}+P_k \gamma_{k+2} x_{k+2}\bigr) \\
    &\qquad - \bigl(\gamma_{k+1}(-1)^k\gamma_1\dots\gamma_{k}\gamma_\ell x_{\ell-1}
        + \gamma_{k+1} P_k x_k+\gamma_{k+1}^2 P_{k-1}x_{k+1}\bigr) \\
    &= (-1)^{k+1}\gamma_1\dots\gamma_{k+1}\gamma_\ell x_{\ell-1}+
    (P_k -\gamma_{k+1}^2 P_{k-1})x_{k+1}+\gamma_{k+2} P_{k}x_{k+2}
\;=\; 0\,.
\end{align*}
Defining $P_{k+1} \coloneqq P_k - \gamma_{k+1}^2 P_{k-1}$ for $k\leq \ell-2$, we conclude that $(L'_{k+1})$
also holds. Since we know that $(L'_1)$ holds, it follows by induction that
$(L'_{\ell-2})$ holds, implying that
\begin{align*}
0 \;=\; (-1)^{\ell-1}\gamma_1\dots\gamma_\ell x_{\ell-1}+
    P_{\ell-1} x_{\ell-1}+\gamma_{\ell} P_{\ell-2}x_{\ell}\,,
\end{align*}
that is,
\begin{align*}
\gamma_{\ell} P_{\ell-2}x_{\ell}
\;=\; \bigl((-1)^{\ell}\gamma_1\dots\gamma_\ell - P_{\ell-1}\bigr) x_{\ell-1}\,.
\end{align*}
To finish the proof, we need to control the sign of the coefficients appearing
above. Once again, the strategy of expressing relevant quantities using the
independent events $A_k$ plays a role.

\begin{lemma}
\label{lema:seq_Pk}
Define 
$$P_0 \coloneqq 1,\quad  P_1\coloneqq 1 - \gamma^2_1,\quad  P_k \coloneqq P_{k-1} - \gamma_k^2
P_{k-2}, \; 2\leq k\leq \ell-1,$$
%
and
%
$$P_\ell\coloneqq 2(-1)^\ell \gamma_1\cdots \gamma_\ell.
$$ 
%
Then
\begin{equation}
\label{eq:seq_Pk}
1    \;=\; P_0 \;>\; P_1 \;>\; P_2 \;>\; \cdots \;>\; P_{\ell-1} > P_{\ell}
    \;=\; 2(-1)^{\ell} \gamma_1 \cdots \gamma_\ell
    \;>\;0\,.
\end{equation}
\end{lemma}

\begin{proof}
Recall $A_k$ as defined in~\eqref{eq:def_event_Ak} and $B_k = A_k \cap
A_{k+1}^c$. We simply notice that the sequence $P_k$ in the statement can be
alternatively described by the equation
\begin{equation}\label{eq:seq_Pk_as_union}
    P_k \;=\; 1 - \bb P \Bigl(\bigcup_{j=1}^k B_j\Bigr),\quad k\leq \ell-1.
\end{equation}
Indeed, it is straightforward to check~\eqref{eq:seq_Pk_as_union} for $k=0,1$.
Now, suppose that~\eqref{eq:seq_Pk_as_union} holds for $k-1$. Then
\begin{align*}
1 - \bb P\Bigl( \bigcup_{j=1}^k B_j\Bigr)
    \;=\; 1 - \bb P\Bigl( \bigcup_{j=1}^{k-1} B_j\Bigr) - \bb P (B_k)
    + \bb P\Bigl( B_k \cap \bigcup_{j=1}^{k-1} B_j\Bigr).
\end{align*}
Since $B_k \cap B_{k-1} = \varnothing$, the intersection above is given by
\begin{equation*}
\bb P\Bigl( B_k \cap \bigcup_{j=1}^{k-1} B_j\Bigr)
\;=\; \bb P\Bigl( B_k \cap \bigcup_{j=1}^{k-2} B_j\Bigr)
    \;=\; \bb P(B_k) \bb P\Bigl(\bigcup_{j=1}^{k-2} B_j\Bigr)
    \;=\; \gamma_k^2 (1 - P_{k-2})\,,
\end{equation*}
%
where in the second identity we used that $B_k=A_k\cap A_{k+1}^c$ and $\cup_{j\leq k-1} B_j=\cup_{j\leq k-2}(A_j\cap A_{j+1}^c)$ are mutually independent, because the $A_j$'s are.
Putting together the equations above, we obtain
\begin{equation*}
1 - \bb P\Bigl( \bigcup_{j=1}^kB_j\Bigr)
    \;=\; P_{k-1} - \gamma_k^2 + \gamma_k^2 (1-P_{k-2})
    \;=\; P_{k-1} - \gamma_k^2 P_{k-2}
    \;=\; P_k\,,
\end{equation*}
completing the induction step.
The inequalities in~\eqref{eq:seq_Pk} are now evident from the fact that
the sequence of events $\cup_{j=1}^k B_j$ is increasing in $k$ and that
we already know $P_{\ell} = \bb P(\cap B_j^c) = 2 (-1)^{\ell} \gamma_1 \dots
    \gamma_{\ell} > 0$, see~\eqref{eq:prob_cup_Bk_2}.
\end{proof}

With Lemma~\ref{lema:seq_Pk}, we can finish the proof by noticing that
\begin{equation*}
\gamma_\ell P_{\ell-2} \;<\;0
\quad \text{and} \quad
(-1)^{\ell}\gamma_1\dots\gamma_\ell - P_{\ell-1}
\;=\; \frac{1}{2}P_{\ell} - P_{\ell-1} \;<\; 0\,,
\end{equation*}
which imply that $x_{\ell-1}$ and $x_{\ell}$ have the same sign. The reasoning
above actually shows that any eigenvector of zero with some positive entry is in $(0,\infty)^\ell$. Finally, we argue that the eigenspace of zero
has dimension 1. The Spectral Theorem ensures that $\Sigma$ has
an orthonormal basis of eigenvectors. Now, suppose that $v_1, v_2$ are two
orthogonal eigenvectors of zero. Replacing $v_j$ by $-v_j$ if needed, we can
assume $v_j \in (0, \infty)^\ell$ for $j=1,2$. But then their inner product is
positive, leading to a contradiction.
\end{proof}

\end{document}
